% Reconstructed from the independently audited final proof after the workspace reset.
\paragraph{C.3aq (Refuted): the integral top quotient for $p>0$.}
For every algebraically closed field $k$ of characteristic $p>0$,
there is a specified excellent normal strictly henselian local
$k$-algebra $R$ of dimension three, with residue field $k$ and
regular puncture $X$, for which there is a surjection
\[
 \mathrm{CH}^2(X)/\mathrm{CH}^2(X)_{\mathrm{alg}}
                                      \longrightarrow(k,+).
\]
Algebraic equivalence uses actual families on fixed products.
Thus integral top finiteness fails, including over every
$\overline{\mathbb F}_p$. The exhibited quotient has exponent
$p$; no rational or characteristic-zero conclusion is asserted.

\begin{proof}
\emph{1. The relative quotient in every characteristic.}
For a split commutative square-zero extension $A\oplus I$,
there is a natural surjection
\[
 \tau:K_2(A\oplus I,I)\longrightarrow
                                  I\otimes_A\Omega^1_{A/\mathbb Z}.
                                                        \tag{C.3aq.1}
\]
The square-zero ideal lies in the Jacobson radical, so the
relative Dennis--Stein presentation in
\citet[Chapter~III, Theorem~5.11.1(b), p.~44]{frWeibelKbookIII}
applies. For $m\in I$ and $r\in A\oplus I$ set
\[
 \tau\langle m,r\rangle=-m\otimes d\bar r,
 \qquad \tau\langle r,m\rangle=m\otimes d\bar r.
                                                        \tag{C.3aq.2}
\]
Both formulas are zero when both entries lie in $I$.
Relation (D1) is skew-symmetry. In (D2), if the first
entry lies in $I$, the reduction of $s+t-rst$ is
$\bar s+\bar t$; if the other two entries lie in $I$,
their product is zero. This gives additivity in both
cases. In (D3),
\[
 \langle r,st\rangle=\langle rs,t\rangle+\langle tr,s\rangle,
\]
the case $r\in I$ is the Leibniz rule. If $s\in I$,
the terms containing $\bar r\,s\,d\bar t$ cancel,
leaving $s\bar t\,d\bar r$. The case $t\in I$
is the same computation with exchanged factors.
These are all relative relations, and the generators
map onto the differential module. This proves (C.3aq.1).

The map is natural in split pairs and linear under
rescaling of the square-zero ideal. Its sign is the
negative of Exercise~5.14(a), in the modern convention
of Definition~5.11 in the same source. Thus
\[
                    \tau\{1+tc,u\}=tc\,d\log u.
                                                        \tag{C.3aq.3}
\]
Indeed $\langle u,-tc/u\rangle=\{u,1+tc\}$ and
skew-symmetry gives the formula. Symbols with two
infinitesimal units have zero image, without being
asserted zero in characteristic two.

For $V_2=\operatorname{Spec}_V(\mathcal O_V\oplus L)$,
the map glues to a quotient of the relative $\mathcal K_2$
sheaf onto $\Omega^1_{V/\mathbb Z}\otimes L$.
The ring splitting provides projection onto the relative
part. Since $k$ is perfect, the target is exactly
$\Omega^1_{V/k}\otimes L$. No isomorphism in characteristic
two or surjectivity on cohomology is assumed.

\emph{2. Uniform explicit geometry.}
Put
\[
 e=p+1,\qquad n=pe=p(p+1),\qquad d=n-3>0,
                         \qquad D_1=ed\,\infty.
\]
In coordinates $(Y:Z:W)$ take
\[
 C:\quad Y^n-YW^{n-1}=Z^{n-1}W,
                       \qquad \infty=(0:1:0).
                                                        \tag{C.3aq.4}
\]
The affine equation $y^n-y=z^{n-1}$ has $y$ derivative
$-1$. At the unique infinity point the $W$ derivative
is $-Z^{n-1}$. Thus $C$ is smooth, $z$ is an etale
affine coordinate, and the genus formula gives
$\operatorname{div}(dz)=n(n-3)\infty=pD_1$.
This is the Tango family in
\citet[Example~1.3 and Lemma~1.4, author pp.~3--4]{D_Mukai2011}.

An exact normalized infinity chart is $q=Y/Z$, $h=W/Z$:
\[
 h=q^n-qh^{n-1},\qquad
 z-q^{-n}=q^{n(n-3)}c,\qquad
 c=q(h/q^n)^{n-2}.
                                                        \tag{C.3aq.5}
\]
The function $q$ is a uniformizer, $h/q^n$ is a unit
with residue one, and $c$ is a uniformizer with $dc$
a generator.

Take the nonzero Frobenius-killed extension
$0\to\mathcal O_C(-D_1)\to E\to\mathcal O_C\to0$
defined by the normalized exact differential $dz$.
The negative line bundles in the Frobenius sequence
have no sections. Hence the class is nonzero and its
Frobenius pullback has a unique splitting.
In the bundle of lines $P=\mathbb P_C(E)$ the infinity
section has normal bundle $\mathcal O_C(D_1)$.
The complementary Frobenius section gives a smooth
integral divisor $G$, disjoint from $S_\infty$, with
affine equation $s^p=z$. Normalization (C.3aq.5)
verifies its smoothness at infinity. As in C.3an,
fibre degree and restriction determine the correct sign:
\[
                   G\sim pS_\infty-p\bar g^*D_1.
                                                        \tag{C.3aq.6}
\]
The explicit source-sign audit in C.3an applies here.

Form the cyclic cover of degree $e$ with inverse root
\[
 M=\mathcal O_P(-S_\infty)\otimes
                          \bar g^*\mathcal O_C(pd\infty),
 \qquad V=\operatorname{Spec}_P\left(\bigoplus_{j=0}^{e-1}M^j\right).
                                                        \tag{C.3aq.7}
\]
Multiplication is defined by the section of $M^{-e}$
with divisor $G+S_\infty$, with scalar fixed by
$T^e=s^p-z$. Formula (C.3aq.6) proves the root identity.
The simple branch valuation proves connectedness, even
if $e$ is composite. The cover is smooth because $e$
is prime to $p$ and the smooth branch components are
disjoint. It is projective because it is finite over $P$.
For $g:V\to C$ and the reduced infinity section $S$,
\[
 N_{S/V}=M^{-1}|_{S_\infty}=\mathcal O_C(d\infty),\qquad S^2=d.
\]
Set
\[
                   D=pS+d\,g^*\infty,\qquad L=\mathcal O_V(pD).
                                                        \tag{C.3aq.8}
\]
Every fibre is the integral compactification of
$T^{p+1}=(s-c_0)^p$ and meets $S$ once. Thus
\[
 D^2=(e^2-1)d>0,\qquad D\cdot S=ed>0,\qquad D\cdot F=p>0.
\]
Every other horizontal curve has nonnegative intersection
with $S$ and positive degree over $C$. These cases cover
every integral curve, so Nakai--Moishezon proves ampleness.

\emph{3. The cotangent section, including $p=2$.}
At infinity normalize by
$s=q^{ed}s_i+q^{-e}$ and $T=q^{pd}T_i$.
Then $T_i^{p+1}=s_i^p-c$. On a finite fibre chart
$s_i,T_i$ are smooth coordinates. At the infinity
section put $x=1/s_i$ and $w=T_i/s_i$; directly,
\[
             w^{p+1}=x-cx^{p+1},\qquad
             dT_i=-\frac{x^p}{w^p}\,dc.
                                                        \tag{C.3aq.9}
\]
The derivative identity uses $T_i=w/x$ and divides
neither by $p$ nor by two. Since $x/w^{p+1}$ is a
unit, its order along $S$ is $p(p+1)-p=p^2$.
Also $dT=q^{pd}dT_i$, since the scalar $pd$ is zero
in $k$. The exact zero divisor is therefore
$p^2S+pd\,g^*\infty=pD$, with no residual zero at
the intersection. This gives a nowhere vanishing section
\[
                 \eta=dT\in H^0(V,\Omega^1_{V/k}\otimes L^{-1}).
                                                        \tag{C.3aq.10}
\]
At $v:y=z=s=T=0$, its normalized value is $dT$
and $s,T$ are etale coordinates.

\emph{4. The actual local ring and integral invariant.}
Define, with residue field $k$,
\[
 A=\bigoplus_{j\ge0}H^0(V,L^j),\quad
 R=(A_{A_{>0}})^{\mathrm{sh}},\quad
 X=\operatorname{Spec}R\setminus\{\mathfrak m_R\}.
                                                        \tag{C.3aq.11}
\]
The full section ring is normal and finitely generated.
The Veronese and consecutive-powers proof in C.3an,
Step~2, gives the proper contraction of
$Z_0=\operatorname{Tot}(L^{-1})$, an isomorphism off
the zero section, without requiring degree-one generation.
The base change $Z$ to $R$ is regular, with complement
$X$, exceptional divisor $V$, and split first neighborhood
$V_2$ with conormal $L$. The local ring has dimension three
and every property in the statement.

The same integral Bloch formula and finite-stage passage
give $\mathrm{CH}^2(Z)=H^2(Z,\mathcal K_2)$.
Restrict to $V_2$, project onto its relative part, and
apply (C.3aq.1), obtaining
\[
 \delta:\mathrm{CH}^2(Z)\longrightarrow
                               H^2(V,\Omega^1_{V/k}\otimes L).
\]
No injectivity of (C.3aq.1) or surjectivity on $H^2$
is used. The integral codimension-two descent edge in
C.3an has multiplier $(2-1)!=1$ in every characteristic.
For an exceptional class $i_*\beta$ with
$\beta\in\mathrm{CH}^1(V)$, its $K_0$ restriction is
the pulled-back class $\pi^*\beta(1-[L])$. Its relative
component vanishes before $\tau$ is applied. Thus
$\delta$ descends to $\mathrm{CH}^2(X)$.

Pair with (C.3aq.10) and the functorial Serre trace:
\[
 \rho:\mathrm{CH}^2(X)\longrightarrow k,
                   \rho(\xi)=\operatorname{Tr}(\delta\xi\wedge\eta).
                                                        \tag{C.3aq.12}
\]
The actual ordinary algebraically trivial subgroup is
divisible, by the proper-curve correspondence proof in
C.3an, Step~4, which is valid in every characteristic.
Its image in the exponent-$p$ target is zero, so $\rho$
factors through the ordinary quotient.

\emph{5. Actual curves attain every scalar.}
On the affine chart of $Z_0$ with fibre coordinate $t$,
fix $a\in k$ and impose
\[
 s=0,\quad T=at,\quad z=-a^{p+1}t^{p+1},\quad
                    y^n-y=(-a^{p+1}t^{p+1})^{n-1}.
                                                        \tag{C.3aq.13}
\]
The last equation has $y$ derivative $-1$. Its specified
point $y=t=0$ is smooth and belongs to a unique
irreducible component, whose germ is transverse to $V$.
For $a=0$ take the component $y=0$. No global
irreducibility assertion about this polynomial is needed.
Close this component, normalize its affine image curve,
and select the finite henselian branch with closed point
$v$, exactly as in C.3an, Step~5. Smoothness at that
germ and properness over the henselian local branch
exclude every other non-isomorphism point. Thus this
gives an actual closed integral curve $D_a\subset Z$
meeting $V$ only at $v$, not a formal curve declared
to be algebraized.

Its derived intersection with $V_2$ is the Koszul
point class of $(s,T-at)$. Relative to $D_0$, formula
(C.3aq.3) gives
\[
 \tau\bigl(\{s,T-at\}-\{s,T\}\bigr)
                           =t\,\frac{a\,ds}{sT}.
\]
This is still true in characteristic two: extra
relative symbols have zero image under $\tau$.
Wedging with $t^{-1}dT$ and taking the ordered
local trace gives $a$. Functorial Serre duality gives
the same global trace. Hence
$\rho([D_a]-[D_0])=a$ for every scalar, proving
surjectivity.
\end{proof}

This extension uses a relative quotient, not an
incorrect characteristic-two tangent isomorphism.
It concerns integral top finiteness in local dimension
three. Rational and complex top quotients, and the
higher-dimensional top cases, are not consequences
of this statement.
